% !TEX program = xelatex
% 问题一：不同风速下的锚链静力形状图。
% 数据由 src/q1/q1_static_equilibrium.py 生成；可从仓库根目录、paper/ 或本目录编译。
% 例如：cd src/q1 && xelatex q1_chain_shapes.tex
\documentclass[tikz,border=5mm]{standalone}
\usepackage[UTF8,fontset=fandol]{ctex}
\usepackage{pgfplots}
\pgfplotsset{compat=1.18}

% 以相对路径定位脚本生成的两种风速下的链坐标。
\IfFileExists{../../data/processed/q1_chain_shape_12ms.csv}{%
  \def\QOneProcessedDir{../../data/processed}%
}{%
  \IfFileExists{../data/processed/q1_chain_shape_12ms.csv}{%
    \def\QOneProcessedDir{../data/processed}%
  }{%
    \def\QOneProcessedDir{data/processed}%
  }%
}

\begin{document}
\begin{tikzpicture}
\begin{axis}[
  width=14cm,
  height=12.6cm,
  title={不同风速下锚链的静力平衡形状},
  xlabel={相对锚点的水平坐标（m）},
  ylabel={高程（m）},
  xmin=-0.8, xmax=18.8,
  ymin=-18.8, ymax=1.2,
  xtick={0,3,6,9,12,15,18},
  ytick={-18,-15,-12,-9,-6,-3,0},
  grid=major,
  grid style={gray!25},
  axis line style={black!70},
  tick style={black!70},
  legend style={
    at={(axis cs:1.2,-2.2)}, anchor=north west,
    draw=black!35, fill=white, fill opacity=0.94, text opacity=1,
    font=\scriptsize
  },
  legend cell align={left},
  axis equal image,
  every axis plot/.append style={line cap=round}
]
  \addplot+[very thick, blue!70!black, no marks]
    table [x=x_m, y=z_m, col sep=comma]
    {\QOneProcessedDir/q1_chain_shape_12ms.csv};
  \addlegendentry{风速 $12\,\mathrm{m/s}$}

  \addplot+[very thick, red!70!black, no marks]
    table [x=x_m, y=z_m, col sep=comma]
    {\QOneProcessedDir/q1_chain_shape_24ms.csv};
  \addlegendentry{风速 $24\,\mathrm{m/s}$}

  \addplot+[blue!55!black, line width=0.9pt, domain=-0.8:18.8, samples=2] {0};
  \addlegendentry{静水面}

  \addplot+[brown!70!black, line width=0.9pt, domain=-0.8:18.8, samples=2] {-18};
  \addlegendentry{海床（$-18\,\mathrm{m}$）}

  \addplot+[only marks, mark=*, mark size=2.2pt, black]
    coordinates {(0,-18)};
  \addlegendentry{锚点}

  \node[font=\scriptsize, fill=white, fill opacity=0.9, text opacity=1,
        inner sep=1.3pt, anchor=west]
    at (axis cs:0.45,-1.0) {锚点为坐标原点，$z$ 轴向上};
\end{axis}
\end{tikzpicture}
\end{document}
